\chapter{Learning Analytics and the OULAD Problem}
\label{ch:learning-analytics}

\begin{learningobjectives}
After completing this chapter, you will be able to:
\begin{itemize}
  \item define learning analytics, engagement, assessment performance, and
    outcome forecasting in plain and technical language;
  \item distinguish descriptive, predictive, and causal questions;
  \item explain the benefits and risks of using educational records;
  \item state what the implemented OULAD project can and cannot establish; and
  \item write a responsible analytical question with an explicit unit and
    claim boundary.
\end{itemize}
\end{learningobjectives}

\section{Vocabulary for asking questions with data}

Chapter 1 introduced the software words. This chapter introduces the words used
to turn a real-world concern into a question that data can answer.

\begin{longtable}{@{}L{0.22\textwidth}L{0.71\textwidth}@{}}
\caption{Basic analytical vocabulary}\label{tab:analytical-vocabulary}\\
\toprule
Term & Meaning in this course \\
\midrule
\endfirsthead
\toprule
Term & Meaning in this course \\
\midrule
\endhead
\term{Record or observation} &
One recorded case. Depending on the table, it might be a course offering, a
student's course attempt, an assessment, or one day's online activity. \\
\term{Variable} &
One recorded characteristic that can differ across observations, such as final
result, score, or number of clicks. In a table, a variable is usually stored in
a column. \\
\term{Measure or metric} &
A value chosen to summarize something, such as attempt count or average score.
A metric is a definition made by an analyst; it is not the complete reality
behind the number. \\
\term{Outcome} &
The event or result being described or predicted. In one OULAD question it
might be the final-result category; in another it might be whether the next
assessment is missing. \\
\term{Population} &
The full set of observations a claim is about. ``All attempts recorded in
OULAD'' is a population; ``all university students everywhere'' is not the
population represented by this dataset. \\
\term{Unit of observation} &
What one record represents. This is also called the observation grain. A
student and a student-course attempt are different units. \\
\term{Proportion and percentage} &
A proportion is a part divided by a whole, such as withdrawn attempts divided
by all attempts. Multiplying a proportion by 100 expresses it as a percentage.
The denominator, or whole, must always be named. \\
\term{Pattern or association} &
A difference or relationship observed in records. An association can be real
without showing that one variable caused the other. \\
\term{Proxy} &
A recorded stand-in for something that is difficult to measure directly.
Clicks can be a proxy for one aspect of online activity, but they are not a
direct measure of attention, effort, or learning. \\
\term{Prediction or forecast} &
An estimate of a later or unknown recorded outcome. A prediction says what the
model expects, not why the outcome will occur. \\
\term{Model} &
A mathematical rule that uses input variables to calculate a summary or
prediction. Models can be useful and still be uncertain or wrong. \\
\term{Claim boundary} &
A sentence stating how far the evidence reaches. It names what a result
describes and what it does not prove. \\
\bottomrule
\end{longtable}

\begin{plainlanguage}
An analytical question should identify the people or records under discussion,
what one record represents, what is measured, and how strong an answer the data
can support. Those four choices prevent a number from quietly becoming a much
larger claim.
\end{plainlanguage}

\section{What learning analytics means}

\term{Learning analytics} is the organized use of data about learners and
learning environments to understand and improve educational processes. The
phrase covers many tasks: describing participation, finding data-quality
problems, evaluating assessment patterns, forecasting a recorded event, or
supporting a conversation about course design.

\begin{plainlanguage}
Learning analytics asks what can be learned from records created while
education is taking place. The records are traces of a system, not complete
measurements of a person.
\end{plainlanguage}

The OULAD project studies records from selected Open University module
presentations, meaning particular deliveries of a course, in 2013 and 2014
\citep{oulad_dataset,oulad_paper}. Its source
tables connect module offerings, registrations, assessments, recorded student
information, online resources, submissions, and Virtual Learning Environment
(VLE) interactions. A VLE is the institution's online course system.

\section{Four concepts that must not be collapsed}

\begin{description}
  \item[Student engagement] A broad educational idea concerning participation
    and involvement. In this project, it is represented only by limited
    indicators such as recorded VLE clicks and active days.
  \item[Assessment performance] Recorded evidence associated with coursework or
    examinations, including submission timing, missing submissions, scores, and
    weights.
  \item[Course outcome] OULAD's recorded final category: Distinction, Pass, Fail,
    or Withdrawn.
  \item[Outcome forecast] A model estimate about a later recorded event using
    information available at a specified earlier point.
\end{description}

The first three concepts describe different objects. A student can click often
without understanding the material; a low click count can coexist with offline
study; a missing submission can reflect circumstances not captured in the
data; withdrawal is an administrative outcome, not a diagnosis.

\section{Description, prediction, and causation}

\begin{table}[htbp]
\centering
\caption{Three kinds of analytical claim}
\label{tab:claim-types}
\begin{tabularx}{\textwidth}{L{0.19\textwidth}L{0.30\textwidth}YY}
\toprule
Claim type & Example OULAD question & What it can establish & What it cannot establish \\
\midrule
Descriptive & How many attempts ended in withdrawal? & A count or proportion in the observed records & Why withdrawal occurred \\
Predictive & Can information available this week rank later assessment risk? & How well a forecast works when checked on records not used to build it & Causal explanation or individual need \\
Causal & Would a new support program reduce withdrawal? & The change caused by the program, if the study fairly compares what would happen with and without it & Not available from these existing records alone \\
\bottomrule
\end{tabularx}
\end{table}

\begin{projectconnection}
The project's saved outcome table records 3,024 Distinction, 12,361 Pass, 7,052
Fail, and 10,156 Withdrawn student-module attempts. The withdrawal proportion
is 31.2 percent. This is a descriptive fact about recorded attempts. It does
not identify a cause.
\end{projectconnection}

\begin{figure}[htbp]
\centering
\includegraphics[width=0.86\textwidth]{outcome_distribution.png}
\caption{Recorded final outcomes across 32,593 student-module attempts. This
verified chart describes the dataset; it does not explain the outcomes.}
\label{fig:outcome-distribution}
\end{figure}

\section{The implemented project questions}

The project moves from description toward prediction through five questions:

\begin{enumerate}
  \item How do engagement volume, consistency, and inactivity differ across
    weeks and module presentations?
  \item How do assessment availability, submission timing, missing submissions,
    and weighted progress, where assessments contribute according to their
    course weights, vary?
  \item What activity patterns appear before recorded withdrawal?
  \item At the end of a course week, how well can available information forecast
    a missing or below-pass next assessment?
  \item Do predicted percentages agree with observed percentages
    (\term{calibration})? How many cases are flagged at a chosen cutoff
    (\term{threshold workload})? Does performance change by course week or on a
    presentation excluded while building the model?
\end{enumerate}

Only the first part of this progression is taught in Edition 0.1. The complete
questions are introduced now so that database design is connected to later
analytical purpose.

\section{Potential value}

Responsible educational analytics can:
\begin{itemize}
  \item reveal where records are incomplete or inconsistent;
  \item compare patterns across courses without assuming they have one cause;
  \item make operational definitions explicit;
  \item support investigation of course design and assessment timing;
  \item expose uncertainty and tradeoffs before a decision; and
  \item help institutions ask better questions about support.
\end{itemize}

These benefits depend on careful data definitions. A count at the wrong grain
can produce a convincing but false conclusion.

\section{Risks and boundaries}

Educational data concern people in settings with unequal resources and power.
Important risks include revealing private information; measuring some students
or behaviors less accurately than others; repeating inequalities already
present in historical records; attaching harmful labels; allowing a prediction
to shape the very outcome it predicts; and using analysis to punish rather than
support.

\begin{validationbox}
Before presenting an educational result, answer:
\begin{enumerate}
  \item What exactly was measured?
  \item At what grain was it measured?
  \item What important behavior is absent from the data?
  \item Is the claim descriptive, predictive, or causal?
  \item What harmful interpretation could a reader make?
  \item What human review and opportunity to question or appeal the result
    would be required before action?
\end{enumerate}
\end{validationbox}

A ``high-risk'' prediction must never be treated as a judgment about ability,
effort, motivation, or character. In the implemented project, the primary
model predicts a later recorded assessment event, not a student's worth or
instructional diagnosis. The secondary withdrawal investigation is explicitly
documented as unsuitable for individual use.

\section{Operational definitions}

An \term{operational definition} states how an abstract idea becomes a
measurable variable. For example:

\begin{itemize}
  \item ``weekly activity'' becomes the sum of \texttt{sum\_click} within a
    course-relative week;
  \item ``active days'' becomes the number of distinct recorded activity days;
  \item ``missing assessment'' becomes an eligible student-assessment pair with
    no matching submission row;
  \item ``withdrawal'' becomes a recorded final result and, where present, an
    unregistration day.
\end{itemize}

An operational definition makes analysis reproducible, but it does not erase
the difference between the variable and the underlying educational concept.

\section{Statistical reasoning preview}

The dataset is not a random sample in which members of a larger population had
a known chance of selection. It covers selected, anonymized modules from one
institution and period. Therefore:

\begin{itemize}
  \item a large row count does not guarantee that a result applies elsewhere;
  \item differences between modules may reflect design, population, policy, or
    measurement;
  \item repeated records from the same student are not separate, unrelated
    people;
  \item associations do not establish causes; and
  \item model performance on later presentations does not establish performance
    at another institution.
\end{itemize}

Formal statistical and modeling treatment is forthcoming. Edition 0.1 uses
these ideas to discipline SQL interpretation.

\section{Knowledge check}

\begin{enumerate}
  \item Why is a VLE click not a direct measure of learning?
  \item Classify each project question as descriptive, predictive, or causal.
  \item What does the recorded withdrawal proportion establish?
  \item Give one potential benefit and two potential harms of educational
    analytics.
  \item What is an operational definition?
\end{enumerate}

\begin{exercisebox}{Exercise 2.1: Repair three overclaims}
Rewrite each claim so that it matches the available evidence:
\begin{enumerate}
  \item ``Students with fewer clicks were less motivated.''
  \item ``Missing work caused withdrawal.''
  \item ``The model identifies students who will fail.''
\end{enumerate}
For each revision, name the measured variable, unit of observation, claim type,
and one missing source of context.
\solutionref{sol:learning-analytics}
\end{exercisebox}

\begin{exercisebox}{Exercise 2.2: Frame an OULAD question}
Write one descriptive OULAD question. It must specify:
\begin{itemize}
  \item the population;
  \item the observation grain;
  \item the variable or event;
  \item the grouping or comparison; and
  \item a statement of what the result would not prove.
\end{itemize}
\solutionref{sol:learning-question}
\end{exercisebox}

\begin{commonmistakes}
\begin{itemize}
  \item Using ``student'' when the row actually represents a
    student-module-presentation attempt.
  \item Describing a platform trace as attention, effort, or understanding.
  \item Treating a forecast as an explanation.
  \item Recommending intervention before defining possible benefits, harms, and
    review.
\end{itemize}
\end{commonmistakes}

\transferheading
For any domain, separate the real-world concept from its recorded proxy. A
retail click is not satisfaction, a medical claim is not health, and an
educational click is not learning. State the proxy, its grain, and its claim
boundary.

\begin{summarybox}
Learning analytics uses educational records to answer descriptive and
predictive questions, but the records are incomplete proxies. OULAD supports a
careful study of registrations, assessments, outcomes, resources, and platform
activity. It does not support causal or diagnostic claims by itself. Every
result needs an operational definition, an observation grain, and a responsible
scope of interpretation.
\end{summarybox}
